\part{Bilan et retours d'expérience}
\label{part:demarche_reflexive}

\chapter{Limites observées en production}
\label{chap:reponse_problematique}

\section{Bilan après 38 mois}

\noindent Après 38 mois sur Metis, plusieurs briques ont dû être reprises après mise en production. Les points suivants sont ceux qui ont le plus dévié du design initial.

\section{Compromis forcés en production}

\subsection{Homonymes}

Sur un dossier familial, père et fils \texttt{DUPONT/JEAN} ont reçu tous les bagages sur le premier passager. L'appairage s'appuyait sur la clé textuelle normalisée. La compagnie a rejeté à l'émission. On a ajouté l'index de position (\texttt{travelerIndex}) et un warning UI, mais le cas reste fragile si l'ordre de saisie change.

\subsection{Segments passifs vs ERP agences}

Le découplage POS en base a éliminé les ADM pour double écriture Sabre/Amadeus sur plus de deux ans. Côté agences, Gestour, IEnterprise et TravelCloud ont des imports très rigides : certains refusent le PNR passif si la ligne comptable ne colle pas à leur convention, ou si la devise ne matche pas leur passerelle.

Nous avons ajouté une table d'exceptions / règles par agence en PostgreSQL. Les comptables peuvent importer ; l'équipe doit mettre à jour la table quand un ERP change de format.

\subsection{\texttt{CanonicalFlightOffer} et les Branded Fares}

L'idée du modèle pivot était de transformer n'importe quel SOAP Sabre ou JSON NDC en une structure commune. Pour environ 90\,\% des vols simples (aller-retour économique), ça tient. Dès Turkish Airlines ou Lufthansa Group et leurs \textit{Branded Fares}, les offres collent au billet des services non dissociables (bagage 23\,kg sur le premier tronçon seulement, sièges limités, conditions d'annulation hybrides).

Étendre l'interface TypeScript à chaque particularité alourdissait le modèle. Le champ \texttt{rawCarrierMetadata: Record<string, any>} a servi d'échappatoire pour afficher ces spécificités sans casser le typage du reste.

\subsection{Le 409 ne sauve pas tous les paniers}

L'interception HTTP 409 avec \texttt{PriceChangeConfirmDialog.vue} couvre surtout les écarts de taxes modestes ou le basculement sur la sous-classe immédiatement supérieure. Quand le vol est fermé ou que le saut dépasse le plafond de politique de voyage, l'utilisateur refuse et relance une recherche.

En test, après un refus, Pinia gardait parfois les identifiants de l'ancienne offre. \texttt{useBookingStore().\$reset()} à la fermeture de modale a coupé ce cas.

\chapter{Incidents de production et choix d'architecture}
\label{chap:bilan_critique_technique}

\section{OAuth2 : refresh silencieux et file d'attente}

Entre recherche, choix des vols et saisie des voyageurs, plusieurs minutes s'écoulent. Les access tokens expirent souvent en 15 à 30 minutes. Si l'expiration tombe pendant des appels parallèles (disponibilités, bagages), une implémentation naïve enchaîne des HTTP 401 et déconnecte l'utilisateur en plein checkout.

L'intercepteur Axios (listing \ref{lst:oauth2interceptor}) pose un verrou \texttt{isRefreshing} : au premier 401, un seul refresh part ; les autres requêtes rejoignent \texttt{failedQueue}. Nouveau token disponible : file vidée, headers mis à jour, requêtes rejouées.

\begin{lstlisting}[language=TypeScript, caption={Intercepteur Axios OAuth2}, label={lst:oauth2interceptor}]
import axios, { AxiosInstance, AxiosRequestConfig, AxiosError } from 'axios';

interface QueuedRequest {
  resolve: (token: string) => void;
  reject: (error: any) => void;
}

export function setupAuthInterceptor(apiClient: AxiosInstance, authService: any) {
  let isRefreshing = false;
  let failedQueue: QueuedRequest[] = [];

  const processQueue = (error: any, token: string | null = null) => {
    failedQueue.forEach(prom => {
      if (error) {
        prom.reject(error);
      } else if (token) {
        prom.resolve(token);
      }
    });
    failedQueue = [];
  };

  apiClient.interceptors.response.use(
    response => response,
    async (error: AxiosError) => {
      const originalRequest = error.config as AxiosRequestConfig & { _retry?: boolean };

      if (error.response?.status === 401 && !originalRequest._retry) {
        if (isRefreshing) {
          return new Promise((resolve, reject) => {
            failedQueue.push({ resolve, reject });
          })
            .then(newToken => {
              if (originalRequest.headers) {
                originalRequest.headers['Authorization'] = 'Bearer ' + newToken;
              }
              return apiClient(originalRequest);
            })
            .catch(err => Promise.reject(err));
        }

        originalRequest._retry = true;
        isRefreshing = true;

        try {
          const newToken = await authService.refreshAccessToken();
          isRefreshing = false;
          processQueue(null, newToken);

          if (originalRequest.headers) {
            originalRequest.headers['Authorization'] = 'Bearer ' + newToken;
          }
          return apiClient(originalRequest);
        } catch (refreshError) {
          isRefreshing = false;
          processQueue(refreshError, null);
          authService.handleSessionExpiration();
          return Promise.reject(refreshError);
        }
      }

      return Promise.reject(error);
    }
  );
}
\end{lstlisting}

\section{Incidents rencontrés en production}

\subsection{Écart de centime sur la taxe IZ (taxe Chirac)}

Le symptôme était net : \texttt{ERR\_PRICE\_MISMATCH} sur l'API NDC Air France, réservations Paris -- New York facturées en USD. Les logs Winston montraient 0,01\,USD d'écart entre le total commande et la somme des taxes ventilées (FR, IZ, QW).

On a d'abord suspecté la devise POS injectée dans \texttt{req.context.posCurrency}. Elle correspondait à celle d'Air France. L'écart venait de l'ordre des opérations : \texttt{Math.round} appliqué taxe par taxe après conversion donnait 124,51\,USD, alors que la conversion globale du total attendu valait 124,50\,USD. Nous avons basculé sur l'arrondi bancaire et ajouté une passe finale dans \texttt{pricingConsistency.utils.js} où la dernière composante de taxe absorbe le centime ($\text{Total} = \text{Base} + \sum \text{Taxes}$).

Le listing \ref{lst:winston_ndc_mismatch} reproduit la trace capturée.

\begin{lstlisting}[language=json, caption={Log Winston --- rejet NDC écart de taxes}, label={lst:winston_ndc_mismatch}]
{
  "timestamp": "2024-04-18T14:22:09.142Z",
  "level": "warn",
  "service": "api-aerial",
  "correlationId": "req-b2b-7f89a1c2",
  "pos": "NYC_0042",
  "action": "ORDER_CREATE_REJECTED",
  "supplier": "AIR_FRANCE_NDC",
  "error": {
    "code": "ERR_PRICE_MISMATCH",
    "httpStatus": 409,
    "expectedTotal": "124.50 USD",
    "calculatedTaxSum": "124.51 USD",
    "breakdown": {
      "baseFare": "850.00 USD",
      "taxFR": "45.20 USD",
      "taxIZ_Solidarity": "18.31 USD",
      "taxQW_Airport": "61.00 USD"
    }
  }
}
\end{lstlisting}

\subsection{Horaires Aegean / Turkish décalés de deux heures}

Alertes agences sur des vols Athènes et Istanbul : selon que le poste client était à Paris ou à Londres, l'heure affichée bougeait. Les XML renvoyaient des dates ISO sans suffixe de fuseau (\texttt{"2024-06-15T18:30:00"}). \texttt{new Date()} côté Node les interprétait en UTC ; le navigateur réappliquait le fuseau local. Parsing via \texttt{date-fns-tz} et table aéroport IATA $\rightarrow$ fuseau IANA.

\subsection{Deadlock InnoDB sous Sequelize}

Pendant une campagne promo chez une agence partenaire, les réservations simultanées ont commencé à échouer avec :
\begin{verbatim}
Deadlock found when trying to get lock; try restarting transaction
\end{verbatim}

Augmenter \texttt{connectionLimit} n'a rien changé : les erreurs restaient sur les écritures de dossier, pas sur les lectures de recherche. Un retry avec \texttt{setTimeout} a fait baisser un peu le taux d'erreur sans supprimer la cause.

La repro est venue d'un script qui lançait plusieurs \texttt{OrderCreate} en parallèle sur le même POS. Les traces InnoDB montraient une inversion d'ordre : une transaction faisait \texttt{UPDATE PointDeVente} puis \texttt{INSERT Dossier}, l'autre l'inverse. Nous avons imposé le même ordre d'update dans Sequelize, passé l'isolation à \texttt{READ COMMITTED}, et posé un \texttt{SELECT ... FOR UPDATE} trié par clé primaire avant la mise à jour des compteurs. Les retries applicatifs sont restés en place comme filet.

\chapter{Travail en équipe et responsabilités en production}
\label{chap:posture_professionnelle}

\section{Organisation du développement}

Les développements suivaient des sprints de deux semaines, avec daily et recettes PO. Sur les sujets sensibles (tarification, écritures passives, OAuth2), les merge requests passaient par une revue systématique sur GitLab ou GitHub. Plusieurs refontes (retarification, unification partielle GDS/NDC) ont nécessité une phase de cartographie des flux avant le codage, faute de quoi les correctifs locaux dans les contrôleurs se cumulaient.

\section{Collaboration Paris / Tunisie}

L'équipe de développement est répartie entre Paris et la Tunisie. Les décisions techniques sont documentées (wiki, commentaires de revue). L'exemple des fuseaux horaires illustre le besoin : un parsing via le constructeur natif \texttt{Date} décalait les horaires sur les postes français ; le correctif commun a consisté à introduire \texttt{date-fns-tz} et une table aéroport $\rightarrow$ fuseau IANA, puis à en faire une règle de revue.

\section{Estimations}

Les premiers chiffrages sur les \textit{ancillaries} NDC sous-estimaient la diversité des règles compagnies. Les tests d'émission avec bébé sans siège ont entraîné une réécriture du mapping. Sur SOAP Sabre, le temps réel a surtout été consommé par la gestion des sessions expirées et des réponses non nominales.

\section{Incidents côté agences}

Lors d'une indisponibilité liée à un certificat intermédiaire expiré sur la passerelle Sabre, le support TourCom a signalé un blocage massif en matinée. Les écarts de taxes ou les doubles segments passifs restent par ailleurs une source d'ADM dont le montant peut atteindre plusieurs milliers d'euros par dossier.
